% !TeX program = xelatex
%%%%%%%%%%%%%%%%%%%%%%%%%%%%%%%%%%%%%%%%%%%%%%%%%%%%%%%%%%%%%%%%%%%%%
% Homework Assignment Article
% LaTeX Template
% Version 1.3.1 (ECL) (08/08/17)
%
% This template has been downloaded from:
% Overleaf
%
% Original author:
% Victor Zimmermann (zimmermann@cl.uni-heidelberg.de)
%
% License:
% CC BY-SA 4.0 (https://creativecommons.org/licenses/by-sa/4.0/)
%
%%%%%%%%%%%%%%%%%%%%%%%%%%%%%%%%%%%%%%%%%%%%%%%%%%%%%%%%%%%%%%%%%%%%%
\documentclass[english, a4paper]{article}

% Language
%%%%%%%%%%%%%%%%%%%%%%%%%%%%
\usepackage[utf8]{inputenc}	% Special characters (i.e. 'á')
\usepackage[LGR, T1]{fontenc}	% Math
\usepackage[english]{babel}		% Syllable-sensitive line skips
\let\latinencoding\relax		% Stupid bug fix

% Word Count
%%%%%%%%%%%%%%%%%%%%%%%%%%%%

\usepackage{verbatim} % For word count input

% Ignore footnotes (optional, disable if required)
%TC:macro \footnote [ignore]

% Count words (requires texcount)

\newcommand{\wordcount}[1]{
  \immediate\write18{texcount -nobib -1 -sum #1 > #1-words.sum}
  \input{#1-words.sum}
}

\newcommand{\wordcountfull}[1]{
  \immediate\write18{texcount -nobib -sum #1 > #1-words.sum}
  \verbatiminput{#1-words.sum}
}

% Cross-referencing
%%%%%%%%%%%%%%%%%%%%%%%%%%%%
\usepackage{hyperref}		% Use hyperlinks in document
\usepackage{crossreftools}

% Cross-referenceable Text
% Usage: \labelText{label}{text}
% When referencing \ref{label} will print "text"
\makeatletter
\def\labelText{%
  \@ifstar{\@labeltextstarred}{\@labeltext}
}
\newcommand{\@labeltextstarred}[2]{%
  \crtcrossreflabel*{#1}[#2]%
}
\newcommand{\@labeltext}[2]{%
  \crtcrossreflabel{#1}[#2]%
}
\makeatother

% Graphs, Tables, etc.
%%%%%%%%%%%%%%%%%%%%%%%%%%%%
\usepackage{graphicx}			% Graphs
\usepackage{array} 				% Arrays

% New columns types L, C , R:
\newcolumntype{L}[1]{>{\raggedright\let\newline\\\arraybackslash\hspace{0pt}}m{#1}}
\newcolumntype{C}[1]{>{\centering\let\newline\\\arraybackslash\hspace{0pt}}m{#1}}
\newcolumntype{R}[1]{>{\raggedleft\let\newline\\\arraybackslash\hspace{0pt}}m{#1}}

% TikZ
\usepackage{tikz}

% Extra Symbol Packages
%%%%%%%%%%%%%%%%%%%%%%%%%%%%
\usepackage{amsmath}
\usepackage{amsthm}
\usepackage{mathrsfs}
\usepackage{amssymb}
\usepackage{amsfonts}
\usepackage{cancel}				% Arithmetic cancellation
\usepackage{relsize}			% \mathlarger and \mathsmaller
\usepackage{mathtools}

% Biblatex
%%%%%%%%%%%%%%%%%%%%%%%%%%%%
\usepackage{csquotes}
\usepackage[style=authoryear, backend=bibtex, doi=false, isbn=false,url=false, eprint=false, natbib]{biblatex}
\addbibresource{ecl.bib} %biblatex only
\AtEveryBibitem{\clearfield{month}}
\AtEveryBibitem{\clearfield{day}}
\usepackage{url}

% Edit Page
%%%%%%%%%%%%%%%%%%%%%%%%%%%%

% Margins
\addtolength{\hoffset}{-2.25cm}
\addtolength{\textwidth}{4.5cm}
\addtolength{\voffset}{-3.25cm}
\addtolength{\textheight}{5cm}
\setlength{\parskip}{6pt}
\setlength{\parindent}{0in}

% Multiple columns
\usepackage{multicol}

% Fancy Headers
\usepackage{fancyhdr} % Headers and footers
\pagestyle{fancy} % All pages have headers and footers
\fancyhead{}\renewcommand{\headrulewidth}{0pt} % Blank out the default header
\fancyfoot[L]{} % Custom footer text
\fancyfoot[C]{} % Custom footer text
\fancyfoot[R]{\thepage} % Custom footer text

% Enables comments in red on margin
\newcommand{\note}[1]{\marginpar{\scriptsize \textcolor{red}{#1}}}

% Lists
%%%%%%%%%%%%%%%%%%%%%%%%%%%%
\usepackage{enumitem}

% For 'labeled' itemize labels
% Usage: \myitem[customitem]\label{mylabel}
\newcounter{myitemcounter}
\makeatletter
\newcommand\myitem[1][]{\item[#1]\refstepcounter{myitemcounter}\def\@currentlabel{#1}}
\makeatother

% For numbered sentences
\newcounter{examplesentences}

% For counting inside theorems
\newcounter{intratheorem}

% Custom Paragraph Alignments
%%%%%%%%%%%%%%%%%%%%%%%%%%%%%%
\def\changemargin#1#2{\list{}{\rightmargin#2\leftmargin#1}\item[]}
\let\endchangemargin=\endlist 

\newcommand{\fixmargin}[1]{\begin{changemargin}{1em}{1em}#1\end{changemargin}}

% For referenceable definitions/statements
% Usage: \ndef{label}{title}{text}
% When referencing: \hyperlink{label}{reftext} will print "reftext"
\newcommand{\cndef}[3]{\fixmargin{\hypertarget{#1}{\textit{\textbf{#2}}.} #3}}

% Referenceable definition/statement
% Usage: \ndef{label}{title}{text}
% When referencing: \ref{label} will print "title"
\newcommand{\ndef}[3]{\fixmargin{\labelText{\textit{\textbf{#2}}}{#1}. #3}}

% Extra
%%%%%%%%%%%%%%%%%%%%%%%%%%%%%%%%%%%%%%%%%%%%%%%%%%%%%%%%%%%%%%%%%%%%%

\usepackage{blindtext} % Package to generate dummy text

\usepackage{float} % Improved interface for floating objects
\usepackage{booktabs} % Enhances quality of tables

\usepackage[yyyymmdd]{datetime} % Uses YEAR-MONTH-DAY format for dates
\renewcommand{\dateseparator}{-} % Sets date separator to '-'

% Custom Symbols
%%%%%%%%%%%%%%%%%%%%%%%%%%%%
% Logic symbols
\newcommand{\thus}{/\kern-3pt\therefore\quad}	% Therefore
\newcommand{\nec}{\square}		% Necessary (required: amssymb)
\newcommand{\pos}{\Diamond}		% Possible (required: amssymb)
%\input{fitch.sty}				% Fitch-style proofs

% Parentheses, Brackets and Quotes
\newcommand{\ang}[1]{\left\langle #1 \right\rangle}				% Angled brackets
\newcommand{\tang}[1]{\langle #1 \rangle}				% Angled brackets in text
\newcommand{\tquote}[1]{\text{`}#1\text{'}}					% Text quotes in math mode
\newcommand{\aquote}[1]{\left\ulcorner #1 \right\urcorner}					% Angled quotes
\newcommand{\taquote}[1]{\ensuremath{\ulcorner}#1\ensuremath{\urcorner}}	% Text angled quotes
\newcommand{\brac}[1]{[\kern-1.7pt[ #1 ]\kern-1.7pt]}			% Double brackets
\newcommand{\tbrac}[1]{\ensuremath{[\kern-1.7pt[}#1\ensuremath{]\kern-1.7pt]}} % Text double brackets

% Math symbols
\newcommand{\comp}[1]{{#1}^\mathsf{c}} % Set complement
\newcommand{\N}{\ensuremath{\mathbb{N}}} % Natural numbers
\newcommand{\Z}{\ensuremath{\mathbb{Z}}} % Integers
\newcommand{\Q}{\ensuremath{\mathbb{Q}}} % Rationals
\newcommand{\R}{\ensuremath{\mathbb{R}}} % Real numbers
\providecommand\C{}\renewcommand{\C}{\ensuremath{\mathbb{C}}} % Complex numbers
\DeclareMathOperator{\dom}{dom} % Domain
\DeclareMathOperator{\cod}{cod} % Codomain

% Fonts
%%%%%%%%%%%%%%%%%%%%%%%
\usepackage{fontspec}
\setmainfont{EB Garamond}

% Document
%%%%%%%%%%%%%%%%%%%%%%%%%%%%%%%%%%%%%%%%%%%%%%%%%%%

\begin{document}

% Title Section
%%%%%%%%%%%%%%%%%%%%%%%%%%%%%%%%%%%%%%%%%%%%%%%%%%%

\title{template_assignment} % Article title
\fancyhead[C]{}
\hrule \medskip % Upper rule
\begin{minipage}{0.24\textwidth} % Left side of title section
\raggedright
Oliver Wolf\\[6pt] % Your lecture or course
oliver.wolf@rutgers.edu
\end{minipage}
\begin{minipage}{0.5\textwidth} % Center of title section
\centering 
\large % Title text size
Skepticism I: Descartes' \textit{Meditations}
\\ % Assignment title and number
\normalsize % Subtitle text size
\emph{} % Subtitle
\end{minipage}
\begin{minipage}{0.24\textwidth} % Right side of title section
\raggedleft
11 September 2026\\[6pt] % Date
Rutgers University % Location
\end{minipage}
\medskip\hrule % Lower rule
\bigskip

% Content
%%%%%%%%%%%%%%%%%%%%%%%%%%%%%%%%%%%%%%%%%%%%%%%%%%%

\section{First Meditation}

\subsection*{The Cartesian Method}

Descartes' worry: what if I don't actually know the things I think I know---for example, that I have a body, that other people exist, that what I perceive is the real world rather than an illusion? Many things I've confidently believed---even things I've had good reasons to believe---have turned out to be false. 

His conclusion: to really know something, you must be absolutely \textit{certain}. In other words, there must be no reason to doubt it. If there is any reason to doubt a belief, that belief is unjustified and does not constitute knowledge.

The resulting methodology: Descartes looks at his most basic and seemingly most secure beliefs, and tries to find reasons to doubt them. If he finds any reason for doubt, he rejects the belief. Whatever withstands this test is knowledge.

\subsection*{The Demolition}

First Target: ordinary perceptual beliefs (e.g., that I'm looking at a tree right now)

\begin{itemize}
	\item Reason for Doubt: Sometimes our senses deceive us. For example, from a distance a mound of dirt might look like a capybara.
	\item Takeaway: Not all of our perceptual beliefs are certain, and so not all of them are knowledge. But this only establishes that we don't know in unusual circumstances, like when we are dealing with objects that are far away. This doesn't provide a good reason to doubt, for example, that I have hands and that I'm sitting in philosophy class right now!
\end{itemize}

Second Target: ordinary perceptual beliefs for real this time (even that I have hands or that I'm sitting at a desk)

\begin{itemize}
	\item Reason for Doubt: I could be dreaming. If I'm dreaming, I could very well be somewhere else entirely, and I might not even have a body.
	\item Takeaway: \textit{None} of our ordinary perceptual beliefs constitute knowledge. But this at least still leaves us with knowledge of abstract, structural facts like those of mathematics. Even if the whole world I seem to perceive is an illusion, it is still the case that $2 + 3 = 5$. How could these beliefs be doubted??
\end{itemize}

Third Target: non-perceptual basic beliefs, like those of mathematics

\begin{itemize}
	\item Reason for Doubt: an Evil Demon could be deceiving me. Not only does the demon create my perceptual illusions, but it also makes me think I am certain when I am not. Perhaps I am making an error when I calculate $2 + 3 = 5$ and the demon simply fools me into certainty that I am right.
	\item Takeaway: we don't even know basic mathematical and logical facts.
\end{itemize}

\section{Second Meditation}

We've found reason to doubt seemingly the most obvious truths. Is there \textit{anything} we can be certain about?

\subsection*{The \emph{Cogito}}

Descartes finds \emph{one} belief that is not subject to doubt:

\begin{itemize}
  \item[] \textbf{\emph{Cogito}:} I exist.
\end{itemize}

Why? Two reasons. One particular to the scenario Descartes is considering, and one general.

\begin{itemize}
  \item[] \textbf{Specific reason:}\\ In the Evil Demon scenario, the Demon has to deceive \emph{me}, so I must exist.\\[-6pt]
  \item[] \textbf{General reason:}\\ To even be able to doubt all my beliefs, as I am doing right now, I must be thinking certain thoughts (that reality is an illusion, etc.). But then if \textit{I'm} thinking, I must exist.
\end{itemize}

Descartes concludes that one knows, if nothing else, that one exists. This cannot be doubted, because someone has to be there to do the doubting.

\section{Skepticism Generally}

Descartes' Dreaming and Evil Demon scenarios are examples of \textit{skeptical scenarios}. These are bizarre scenarios in which everything would seem to us as it does now, but in fact many of our beliefs would be false.

\begin{itemize}
	\item Skeptical scenarios are used to try to show that some large domain of beliefs (e.g. perceptual beliefs) do not constitute knowledge. A \textit{skeptic} is someone who believes that we don't know anything in the relevant domain.
	\begin{itemize}
		\item E.g., you can be a skeptic \textit{about} perceptual beliefs and mathematics (as Descartes was) via a skeptical scenario that threatens both perceptual and mathematical beliefs (e.g., the Evil Demon scenario), but one could alternatively be a skeptic about perceptual beliefs without being a skeptic about mathematics---for example, if you thought that something was wrong with the Evil Demon scenario, but endorsed the Dreaming scenario. 
	\end{itemize}
	\item Descartes' Evil Demon scenario is notable for being perhaps the most extreme skeptical scenario, since it makes you doubt not only all your perceptual beliefs, but also even seemingly the most certain domains of inquiry like mathematics.
	\item But there are other well-known skeptical scenarios that threaten smaller domains of knowledge. Here are a few:
	\begin{itemize}
		\item Solipsism: the physical world I perceive is real, but I am the only conscious being. Other minds do not exist.
		\item The Brain in a Vat: the world we perceive is not real because we are all really just brains hooked up to a computer. The computer sends electrical signals to the brains which cause them to hallucinate the world we perceive.
		\item There is no physical world at all---only our minds. The mind of God creates in us the illusion of a physical world.
	\end{itemize}
\end{itemize}

Crucially, a skeptic is not trying to argue that the skeptical scenario actually obtains. Rather, she argues that we can't be \textit{certain} that the scenario doesn't obtain, and that this is enough to erode our knowledge. To understand this better, we need to learn...

\subsection*{A Little Bit of Epistemology}

Skepticism is a central topic in \textit{epistemology}---the study of knowledge. A well-known standard account of knowledge looks roughly like this:\\

An agent $S$ knows a proposition $p$ if and only if:
\begin{enumerate}
	\item S believes that $p$.
	\item $p$ is true.
	\item $S$ is epistemically justified in believing that $p$.
\end{enumerate}

Skeptics are trying to target condition (3): they argue that their skeptical scenario undermines our \textit{justification} in our beliefs. They claim that since our beliefs \textit{could} be false---namely, if the skeptical scenario obtained---we aren't epistemically justified in believing them.

If the skeptic is right, then even if our beliefs happen to be true, they don't count as knowledge; rather, they are merely accidentally true beliefs. For example, even if in fact I'm \textit{not} dreaming, if Descartes is right then I still don't \textit{know} that I have hands: I merely believe that I have hands, and I happen to be correct.

\subsection*{The Certainty Principle}

How do skeptical scenarios undermine our justification in our beliefs, according to the skeptic? Here, we need to extract a suppressed premise.

For the skeptic's argument to work, she needs to assume something like the following:

\begin{itemize}
	\item[] \textbf{Certainty Principle:}\\[4pt] You know that $P$ only if you can rule out every scenario in which $P$ is false (i.e., in which $\neg P$).
\end{itemize}

Without the Certainty Principle, the possibility of a skeptical scenario does not entail that we lack knowledge.

\section{Formalizing Skeptical Arguments}

Now that we have a better understanding of skeptical arguments and how they work, let's try to formalize a general skeptical argument so we can evaluate it further. Here's my attempt:\\

Where $P$ is some ordinary belief, e.g. that I have hands, ...

\begin{enumerate}[label={(\arabic*)}]
	\item I believe that $P$. (Premise)
	\item There is a skeptical scenario SC such that
	\begin{itemize}
		\item $P$ is false in SC, and
		\item my total evidence is exactly the same as it is now. (Premise)
	\end{itemize}
	\item Therefore, I cannot rule out SC. (From 1, 2)
	\item \textbf{Certainty Principle:} I know that $P$ only if I can rule out all scenarios in which $P$ is false. (Premise)
	\item Therefore, I do not know that $P$. (From 3, 4)
\end{enumerate}

\section{Resisting Skepticism?}

Today, most philosophers think that it's probably impossible to refute the skeptic if you grant all of her assumptions. Rather, skepticism can only be resisted by rejecting one of the skeptic's starting assumptions.

Rejecting the Certainty Principle is one good way to do that. The Certainty Principle is extremely demanding. It implies that:
\begin{itemize}
	\item Knowledge requires eliminating all falsifying alternatives, i.e., all $\neg P$ scenarios.
	\item If even one indistinguishable $\neg P$ scenario remains, knowledge fails.
\end{itemize}

But rejecting the Certainty Principle is not without cost. Though it is demanding, the Certainty Principle also has intuitive appeal: if all of my evidence is compatible with a scenario in which $p$ is false, do I \textit{really} know that $p$ is true??\\

For the rest of Unit 1, we'll be looking at two different kinds of replies to skeptical arguments.

\begin{itemize}
	\item Moore's strategy is similar to those just discussed: he thinks we don't need to accept the premises required to get the skeptic's argument going.
	\item Rinard's strategy is more ambitious: in contrast to most philosophers, Rinard thinks we \textit{can} refute the skeptic on her own terms. 
\end{itemize}

\end{document}
